% =====================================================================
%  Chapter 03 — Architecture
% =====================================================================
\chapter{Architecture: Six Phases and Four Layers}
\label{ch:architecture}

A research programme that intends to run for four years and to be executed in
sessions, by more than one agent, with results that must remain checkable
afterwards, needs two structures: a plan for the mathematics and a plan for
the code. This chapter records both as they were specified, and says what
became of each. The verification discipline that grew around them is
Chapter~\ref{ch:house-rules}.

\section{The six-phase proof architecture}
\label{sec:six-phases}

The plan of record divides the intended proof into six phases, with
overlapping schedules over forty-eight months.

\begin{table}[ht]
\centering
\small
\begin{tabular}{@{}c p{0.16\textwidth} p{0.42\textwidth} l@{}}
\toprule
Phase & Months & Goal & Status \\
\midrule
0 & 1--3 & Local existence for Route~1; continuation criterion: a singularity
  implies \(A(T)\to0\) or \(\mu(T)\to0\) & \OPEN \\[3pt]
1 & 1--8 & The scalar master equation for \(T\) globally well posed;
  \(T\in[T_{\min},T_{\max}]\); \(A(T)\), \(\mu(T)\) uniformly bounded
  & \OPEN \\[3pt]
2 & 4--18 & Four regularisation mechanisms; the self-similar-profile lemma
  strengthened & \OPEN \\[3pt]
3 & 12--28 & Global strong solutions for Route~1; the continuation criterion
  never triggered & \OPEN \\[3pt]
4 & 24--36 & The limit \(\mu(T)\to\nu\); the LPS margin \(\varepsilon\)
  survives; Prodi--Serrin gives the Prize & \OPEN \\[3pt]
5 & 32--48 & Assemble; 3D adversarial experiments; peer review; submit
  & \OPEN \\
\bottomrule
\end{tabular}
\caption[The six-phase proof architecture]{The six-phase architecture as
specified. Every phase is \OPEN. The programme did not execute this plan; the
reason is given in \S\ref{sec:plan-vs-path}.}
\label{tab:six-phases}
\end{table}

Three of the six phases were nevertheless \emph{instrumented}, and the
instruments returned results even though the theorems were never proved.
Phase~0's continuation criterion became
\texttt{layer4/continuation\_criterion.py} and was applied to every
adversarial run; it never triggered. Phase~4's limit became
\texttt{layer2/mu\_limit\_2D.py} and \texttt{layer3/mu\_sweep\_3D.py}, twenty
runs in total, all \textsf{PASS}. Phase~2's self-similar-profile question
became the \(\lambda\) sweep of \S\ref{sec:lambda-sweep-plan} below. What was
never built is any part of Phases~1, 3 or 5.

\section{The layer stack}
\label{sec:layer-stack}

The code is organised as four layers with a strict dependency order, and the
order is enforced by a rule: no Layer-3 experiment may be run before Layers~1
and~2 are validated. The stack is drawn in Figure~\ref{fig:layer-stack}.

\begin{figure}[ht]
\centering
\begin{tikzpicture}[
  >={Latex[length=2.2mm]},
  lay/.style={draw, thick, rounded corners=3pt, align=left, inner sep=5pt,
              font=\small, text width=57mm, minimum height=15mm},
  side/.style={font=\scriptsize\itshape, text=rulegrey, align=left,
               text width=34mm},
  ar/.style={->, thick}
]

\node[lay, fill=panelgrey] (l1) at (0,0)
  {\textbf{Layer 1 --- Property engine}\\
   \texttt{layer1/tfirst\_props.py}\\[2pt]
   \footnotesize \texttt{scalar\_T\_props}, \texttt{route1\_coeffs},
   \texttt{route2\_theta}, \texttt{A\_field}, \texttt{second\_law\_check}};
\node[side, right=5mm of l1] (s1)
  {single source of all thermophysical properties; \(A>0\) asserted on every
   call};

\node[lay, above=9mm of l1] (l2)
  {\textbf{Layer 2 --- 2D spectral solvers}\\
   \texttt{layer2/} \quad (M0--M4)\\[2pt]
   \footnotesize \texttt{T\_solver\_2D}, \texttt{route2\_2D},
   \texttt{self\_similar\_IC}, \texttt{lambda\_sweep\_2D},
   \texttt{mu\_limit\_2D}};
\node[side, right=5mm of l2] (s2)
  {cheap, fully resolved, where mechanisms are first isolated};

\node[lay, above=9mm of l2] (l3)
  {\textbf{Layer 3 --- 3D spectral solvers}\\
   \texttt{layer3/} \quad (M5--M9)\\[2pt]
   \footnotesize \texttt{route1\_3D}, \texttt{route2\_3D} (+JAX ports),
   \texttt{wang\_profiles\_3D}, \texttt{mu\_sweep\_3D},
   \texttt{blowup\_search\_3D}};
\node[side, right=5mm of l3] (s3)
  {the real geometry; \(N=32\) to \(256\); where vortex stretching exists at
   all};

\node[lay, above=9mm of l3] (l4)
  {\textbf{Layer 4 --- Diagnostics}\\
   \texttt{layer4/} \quad (M10 and after)\\[2pt]
   \footnotesize \texttt{alignment\_bridge}, \texttt{k\_correlation},
   \texttt{band\_concentration}, \texttt{scale\_diagnostics},
   \texttt{continuation\_criterion}, \dots};
\node[side, right=5mm of l4] (s4)
  {each module targets one labelled statement in the proof document};

\draw[ar] (l1) -- (l2);
\draw[ar] (l2) -- (l3);
\draw[ar] (l3) -- (l4);

\draw[ar, dashed] (l4.west) to[out=180, in=180, looseness=0.85]
  node[left, font=\scriptsize, align=right, pos=0.5, inner sep=1pt]
  {results\\feed the\\proof doc.} (l1.west);

\node[draw=rulegrey, dashed, rounded corners=3pt, fit=(l1)(l4),
      inner sep=9pt] (stack) {};
\node[below=1mm of stack, font=\scriptsize\itshape, text=rulegrey]
  {strict dependency order: nothing above runs until everything below is
   validated};
\end{tikzpicture}
\caption[The layer stack]{The four-layer stack. The dependency order is a
rule, not a convention: it exists so that a failure at Layer~3 or~4 can be
attributed to the physics rather than to a property lookup or a
two-dimensional bug.}
\label{fig:layer-stack}
\end{figure}

The single architectural rule that did the most work is the one at the bottom:
\emph{all thermophysical property lookups go through Layer~1}, and no solver
file may inline a property formula. It sounds like hygiene. Its real function
is that when \(A_{\min}\) is observed to equal \(8.7236\times10^{-6}\) in an
adversarial 3D run at \(T=150\,\mathrm{K}\), and \(3.0801\times10^{-5}\) in a
Taylor--Green run at \(300\,\mathrm{K}\), those two numbers came from the same
code path and their ratio is meaningful. Chapter~\ref{ch:layer3} uses exactly
that comparison to establish that \(A_{\min}\) is set by the background
temperature and not by the flow --- a conclusion that would have been
unavailable had the two solvers each carried their own copy of the viscosity
law.

\section{Experiment identity}
\label{sec:exp-ids}

Every experiment carries a unique identifier of the form
\texttt{EXP-\{LAYER\}-\{ROUTE\}-\{NNN\}}, optionally with a profile tag:
\texttt{EXP-L2-R1-CCF-003} is the third Layer-2, Route-1 experiment on a
Couette--Couette-flow self-similar profile;
\texttt{EXP-L3-R2-ANISO0100-JAX-128} is a Layer-3, Route-2 anisotropic-collapse
run at aspect ratio \(0.100\) on the JAX backend at \(N=128\).

Every result is logged as
\[
  \texttt{claim\_id} \;\mid\;
  \texttt{verdict}\in\{\textsf{PASS},\textsf{FAIL},\textsf{PARTIAL}\}
  \;\mid\; \texttt{key\_metric} \;\mid\; \text{figure path},
\]
into \texttt{results/results.db}. The database has thirty-one tables. The rule
that gives this force --- \emph{a result without a logged claim\_id is not a
result} --- is discussed in Chapter~\ref{ch:house-rules}, where it is also
recorded that the programme broke it once, in \Sn{30}.

\section{The self-similar-profile interface}
\label{sec:lambda-sweep-plan}

Phase~2's target is the family of self-similar profiles parametrised by a
scaling exponent \(\lambda\), used as adversarial initial data. The interface
is a single module, \texttt{layer2/self\_similar\_IC.py}, which is the sole
source of these profiles for the entire programme; the \(\lambda\) values are
fixed by specification and may not be guessed:
\[
  \begin{array}{llll}
    \text{CCF, stable} & \lambda=1.1808 &
    \text{Boussinesq, stable} & \lambda=1.9206 \\
    \text{CCF, 1st unstable} & \lambda=0.6057 &
    \text{Boussinesq, 1st} & \lambda=1.3991 \\
    \text{CCF, 2nd unstable} & \lambda=0.4703 &
    \text{Boussinesq, 3rd} & \lambda=1.1843
  \end{array}
\]
with the sweep run over \(\lambda\in\{1.2,\,0.9,\,0.6057,\,0.4703,\,0.3,\,0.1\}\)
plus an adversarial minimum at \(\lambda=0.01\). The conjecture under test is
that the viscosity scaling defeats self-similar blow-up for \emph{all}
\(\lambda>0\), with no critical threshold \(\lambda_c\).
Chapter~\ref{ch:layer12} reports what the sweep returned, and how much weight
it can bear.

\section{Plan against path}
\label{sec:plan-vs-path}

It is worth being explicit about the gap between
Table~\ref{tab:six-phases} and what happened, because a reader who takes the
architecture at face value will misread everything after
Chapter~\ref{ch:layer3}.

\begin{figure}[ht]
\centering
\begin{tikzpicture}[
  >={Latex[length=2mm]},
  ph/.style={draw, rounded corners=2pt, font=\scriptsize, minimum width=17mm,
             minimum height=7mm, align=center},
  ev/.style={draw, thick, rounded corners=2pt, font=\scriptsize,
             minimum width=21mm, minimum height=7mm, align=center,
             fill=panelgrey},
  ar/.style={->, thick}
]
\node[font=\small\bfseries] at (-1.6, 1.2) {plan};
\node[ph] (p0) at (0.6,1.2) {Phase 0};
\node[ph] (p1) at (2.9,1.2) {Phase 1};
\node[ph] (p2) at (5.2,1.2) {Phase 2};
\node[ph] (p3) at (7.5,1.2) {Phase 3};
\node[ph] (p4) at (9.8,1.2) {Phase 4};
\node[ph] (p5) at (12.1,1.2) {Phase 5};
\draw[ar, rulegrey] (p0)--(p1); \draw[ar, rulegrey] (p1)--(p2);
\draw[ar, rulegrey] (p2)--(p3); \draw[ar, rulegrey] (p3)--(p4);
\draw[ar, rulegrey] (p4)--(p5);

\node[font=\small\bfseries] at (-1.6,-1.2) {path};
\node[ev] (a) at (0.9,-1.2) {\Sn{1}--\Sn{17}\\numerics};
\node[ev] (b) at (3.9,-1.2) {\Sn{18}--\Sn{27}\\Claim A};
\node[ev] (c) at (6.6,-1.2) {\Sn{29}--\Sn{30}\\Prize claim};
\node[ev] (d) at (9.0,-1.2) {\Sn{31}\\\textsc{audit}};
\node[ev] (e) at (12.0,-1.2) {\Sn{32}--\Sn{47}\\pincer, walls};
\draw[ar] (a)--(b); \draw[ar] (b)--(c); \draw[ar] (c)--(d); \draw[ar] (d)--(e);

\draw[->, dashed, rulegrey] (p2.south) to[out=-90,in=90]
  node[right, font=\scriptsize, pos=0.55] {only the $\lambda$ sweep} (a.north);
\draw[->, dashed, rulegrey] (p4.south) to[out=-90,in=60]
  node[right, font=\scriptsize, pos=0.4] {only the $\mu$ sweep} (a.north east);
\end{tikzpicture}
\caption[Plan against path]{The planned phase sequence (top) against the
programme's actual path (bottom). Only two phase targets --- the
self-similar-profile sweep of Phase~2 and the \(\mu(T)\to\nu\) limit of
Phase~4 --- were reached, and both numerically rather than analytically. From
\Sn{18} the programme worked in Route~2 and the phase plan ceased to describe
it.}
\label{fig:plan-vs-path}
\end{figure}

The reason for the divergence is not that the phase plan was abandoned by
decision. It is that \Sn{18} found a route into the exact Prize equations
that appeared to be closing --- Claim~A, then the Route~C bootstrap, then a
sequence of increasingly specific gap characterisations --- and each session
followed the previous one's most promising thread.
Chapter~\ref{ch:claim-a} records that sequence in detail. By \Sn{29} the
programme believed it had a proof for \(u_0\in H^1(\T^3)\); by \Sn{31} it did
not. Everything after \Sn{31} is a rebuild, and the rebuild took the
architecture's layers with it but not its phases.
